\section{The hemicycle and the main result}\label{rw:main-results}

The quadratic field
\[
 \dot x=-y,\qquad \dot y=x(1+y)
\]
has a center whose period annulus fills the half-plane $y>-1$.
On the Poincar\'e sphere its outer boundary is a hemicycle: the
compactified line $y=-1$ joins two semi-hyperbolic points at infinity,
and the upper equatorial semicircle joins them through a nilpotent
point. We denote this graphic by $\Gamma=H_{14}^{3}$.
The question is how many isolated periodic orbits can appear near
this entire boundary under arbitrary quadratic perturbations.

The reduction of the quadratic finiteness problem to cyclicity
questions for a finite collection of graphics was developed by
Dumortier, Roussarie and Rousseau~\cite{DRR1994}.
For center graphics through a nilpotent point, Roussarie and Rousseau
\cite{RoussarieRousseau2015} prove finite cyclicity of the boundary
limit periodic sets for several neighboring configurations.
Their Theorem~1.1 explicitly leaves $H_{14}^{3}$ aside.
Their family $\dot x=-y+B_0x^2$, $\dot y=x+xy$ is treated for
$B_0>0$; the field above has $B_0=0$.
At the end of their introduction they identify the two
semi-hyperbolic equatorial points as the additional difficulty
in this hemicycle. A bound near a fixed positive $B_0$ does not
give a bound uniform as these points degenerate at $B_0=0$.
The hemicycles studied by Mar\'in and Villadelprat
\cite{MarinVilladelprat2025} have two hyperbolic saddles at infinity;
the semi-hyperbolic and nilpotent points here require additional
slow-passage estimates.

We use the ordinary Euclidean topology on the twelve coefficients
of a planar quadratic vector field. All neighborhoods of $\Gamma$
are taken on the closed northern hemisphere. A cycle is counted only
when its entire affine periodic orbit lies in that neighborhood.
Its multiplicity is the order of the isolated zero of a local
Poincar\'e displacement; a period annulus contributes no isolated
cycles.

\begin{theorem}[Cyclicity three]\label{rw:main:cyclicity}
Let $X_0=(-y,x(1+y))$ and let $\Gamma$ be its hemicycle described
above.
\begin{enumerate}[label=\textup{(\roman*)}]
\item There are an open neighborhood $\mathscr U$ of $\Gamma$ and
an open coefficient neighborhood $\mathscr V$ of $X_0$ such that
every quadratic field in $\mathscr V$ has total multiplicity at most
three among its isolated periodic orbits contained in $\mathscr U$.
\item For every open neighborhood $\mathscr U_1$ of $\Gamma$ and
every open coefficient neighborhood $\mathscr V_1$ of $X_0$, some
quadratic field in $\mathscr V_1$ has three distinct simple limit
cycles contained in $\mathscr U_1$.
These fields can be chosen in a sequence tending to $X_0$ such that
each of the three cycles converges to the entire $\Gamma$ in spherical
Hausdorff distance.
\end{enumerate}
Consequently the cyclicity of $H_{14}^{3}$ in the full quadratic
family is exactly three.
\end{theorem}

The upper bound controls all isolated cycles whose entire orbits lie
in the fixed neighborhood, with their full return multiplicities.
The lower bound also gives cyclicity three in the usual Hausdorff
sense of \cite[Definition~2.1]{RoussarieRousseau2015}. It is realized
in the five-parameter family
\[
 \dot x=-y+Bx^2+my^2+ax,\qquad
 \dot y=x(1+y)+cy+dy^2.
\]
Near-identity affine transformations and positive time changes make
this family a local slice for every quadratic perturbation. The two
parts of the theorem are proved independently in
Theorems~\ref{rw:upper-bound} and~\ref{rw:three-simple-cycles}.

\begin{figure}[htbp]
\centering
\begin{tikzpicture}[x=1cm,y=1cm,>=Stealth,font=\small]
 \draw[very thick] (-3.2,0)--(3.2,0);
 \draw[very thick] (3.2,0) arc[start angle=0,end angle=180,radius=3.2];
 \draw[->,very thick] (-2.15,0)--(-1.55,0);
 \draw[->,very thick] (1.55,0)--(2.15,0);
 \draw[->,very thick] (25:3.2) arc[start angle=25,end angle=40,radius=3.2];
 \draw[->,very thick] (130:3.2) arc[start angle=130,end angle=145,radius=3.2];
 \fill (-3.2,0) circle (2pt) node[below left] {$p_-$};
 \fill (3.2,0) circle (2pt) node[below right] {$p_+$};
 \fill (0,3.2) circle (2pt) node[above] {$q_+$};
 \node at (0,2.55) {upper equator};
 \node[below] at (0.95,-0.12) {$y=-1$};
 \draw[densely dashed,thick] (0,-0.25)--(0,0.72)
       node[above] {$\Sigma$};
 \draw[->,thin] (0.25,0.1)--(0.25,0.48);
 \node[right] at (0.3,0.42) {$z=1+y$};
 \node at (0,1.5) {period annulus side};
\end{tikzpicture}
\caption{The base hemicycle, as a topological schematic of the northern
compactification. Solid curves are the invariant finite line and the
upper equatorial arcs; arrows on these curves give the base field's
time direction. The thin arrow labels the central height, and the
dashed segment is the central transverse section. The vertices
$p_\pm$ are semi-hyperbolic and $q_+$ is nilpotent. Flattening the
compactified finite line to a diameter is schematic; distances and
curvature have no metric meaning.}
\label{rw:fig:hemicycle}
\end{figure}


\subsection{The mechanism of the proof}

Two difficulties intervene between local passage estimates and a sharp
bound. Several passages can lose their exit simultaneously, making
their linear contraction actions arbitrarily large. Near the nilpotent
end, a coefficient direction can also reveal a second scale with
unrestricted separation from the first. These phenomena require
estimates uniform in the complete passing domains. Even with such
estimates, adding the counts of neighboring charts would exceed three.
We instead compare the same two actual paths on one interval and
keep the direction in which their critical values cross zero.

The two paths start at the same outgoing point on $y=0$: one follows
the lower arc backward to the central section, and the other follows
the upper arc and the incoming lower arc forward. We compare their
outputs in the original central height $z=1+y$.
They have outputs $s_++U$ and $s_-+V$, where the reference
difference $s_--s_+$ depends on the field and is independent of the
section input. With $\epsilon>0$ the primary scale of
Definition~\ref{rw:normalization}, we study
\[
 F=U-V-(s_--s_+),\qquad
 K=\epsilon\log(V'/U').
\]
The derivatives $U',V'$ have the same nonzero sign, and
\[
 F'=U'(1-e^{K/\epsilon}).
\]
Hence $F'$ and $K$ have the same zero multiplicities. Once $K$ has
at most two zeros on that interval, Rolle's theorem gives at most
three closing zeros. If a sufficiently small cycle occurs before
the interval, it is unique and simple. It supplies a signed entrance
at which the bound for $K$ improves to one, leaving at most two
further closing zeros. Section~\ref{rw:counting} states the precise
thresholds and proves this counting implication.

Two additional facts make this implication applicable to the whole
hemicycle. The bounded-energy orbits and the large-energy orbits on
which the coefficients remain small are counted together by one
source return. Whenever one monotone comparison does not already
settle the outer cycles, an outer cycle supplies a coefficient
inequality and bounds three equilibrium-free returning strips.
Every intermediate input, including a nonclosing one, then
has both complete central outputs. This is the physical interval
proved in Proposition~\ref{rw:saturation}. The source bound is not
added to a count on this interval when it already includes those
source cycles.

Proposition~\ref{rw:ordered-cover} arranges the physical interval as
a regular prefix followed by at most one terminal union. The regular
crossings have one direction and supply the signs at its entrance.
The hardest endpoint charts initially allow four closing zeros.
Their derivative estimates alone therefore do not prove the theorem.
On a fixed coefficient fibre, Proposition~\ref{rw:coefficient-transport}
transports the outer-cycle inequality to each actual open overlap.
The three-saddle and weak-end arguments then use the already fixed
entrance to obtain a two-zero bound for $K$ on the whole union,
with a one-zero bound under the small-cycle entrance.
These are the exceptional steps in Propositions
\ref{rw:anchors:three-saddle} and \ref{rw:anchors:primary-weak};
Proposition~\ref{rw:anchored-unions} joins them to the other endpoints.

At secondary scale the entrance is carried through the
primary--wedge--secondary regular chain; the actual saddle ratio
and the increasing weak numerator give the terminal bounds in
Propositions~\ref{rw:secondary:saddle-anchor}
and~\ref{rw:secondary:weak-anchor}.

Sections~4--5 construct the original passages and their real center
comparison domains before coefficient division. Relative amplitudes
are estimated separately from possibly divergent contraction actions.
This distinction permits simultaneous small gaps, moving weak roots
and exact cancellation of common powers. Sections~7--10 then classify
the finite, infinite and secondary endpoints on their whole passing
domains. The long secondary passage retains its exact scale exponent
even when the product of the primary scale and the logarithm of the
secondary scale is unbounded.

The lower bound starts from the reversible face $B=0$, $m=\epsilon$.
Section~\ref{rw:lower:preview} describes its three-root variation after
the basic passages and real centers are available.
Section~\ref{rw:lower-bound-section} constructs the coefficient curve,
preserves all three real return paths for one field, and proves that
each resulting cycle approaches the entire hemicycle.

The appendices provide the independent Abelian-integral threshold,
the common-domain calculus, the endpoint algebra and uniform relative
estimates. The domains and derivative orders are stated with the
transitions and endpoint comparisons that use them.
